% ============================================================
% SOLUTION — PART II: MULTIPLE CHOICE
%
% AI INSTRUCTION:
% - Open with the answer table; multi-select answers are listed as "b, d"
% - For each question: restate the correct choice with \correct{...},
%   then explain why it is right AND why each strong distractor is wrong
% - Distractor analysis is the point of an MCQ key. Do not skip it.
% ============================================================

\section*{Part II: Multiple Choice --- Solution}

\begin{center}
\begin{tabular}{|c|c||c|c|}
\hline
\textbf{Q} & \textbf{Answer} & \textbf{Q} & \textbf{Answer} \\
\hline
2.1 & \correct{b} & 2.4 & \correct{a, b, d} \\
2.2 & \correct{c} & 2.5 & \correct{d} \\
2.3 & \correct{b} & & \\
\hline
\end{tabular}
\end{center}

\vspace{4mm}

\textbf{2.1 --- \correct{(b)}}
\begin{ans}
Why (b) is correct, in one or two sentences.

\textbf{Why not the others:}
\begin{itemize}
    \item \textbf{(a)} --- belongs to a different mechanism; name it.
    \item \textbf{(c)} --- true in general, but not what the question asked.
    \item \textbf{(d)} --- a plausible-sounding term that does not exist / does not apply.
\end{itemize}
\end{ans}

\textbf{2.2 --- \correct{(c)}}
\begin{ans}
Model answer with distractor analysis.
\end{ans}

\textbf{2.3 --- \correct{(b)}}
\begin{ans}
Model answer. Show the computation if the question was numeric.
\end{ans}

\begin{keypoint}[Formula]
$$ \text{the formula the question tested} $$
State when it applies and what each symbol means.
\end{keypoint}

\textbf{2.4 --- \correct{(a), (b), (d)}} \quad \textit{(multi-select)}
\begin{ans}
Treat each option as its own true/false judgment and justify all four.
\end{ans}

\textbf{2.5 --- \correct{(d)}}
\begin{ans}
Model answer.
\end{ans}
